\pdfoutput=1
\documentclass[conference]{IEEEtran}
\usepackage{times}
\usepackage[numbers]{natbib}
\usepackage{amsmath,amssymb,booktabs,graphicx}
\usepackage[bookmarks=true,hidelinks]{hyperref}
\hypersetup{pdftitle={Act to Think: Coordinating Information Acquisition and Resource-Aware Reasoning},pdfauthor={Qining Shen}}
\hyphenation{elapsed transitions PPO RARRL Lecture KnowNo}
\newcommand{\E}{\mathbb{E}}
\title{Act to Think: Coordinating Information Acquisition\\and Resource-Aware Reasoning}
\author{\IEEEauthorblockN{Qining Shen}\IEEEauthorblockA{Johns Hopkins University\\601.697 Human-Centered Robotics --- Midterm Report\\November 11, 2026}}
\begin{document}
\maketitle
\begin{abstract}
An embodied agent may lack external evidence or may not have adequately processed the evidence already available. These difficulties call for different choices: acquire information, spend computation on reasoning, or execute a task action. Act to Think studies how to coordinate these choices under resource limits. We formulate the proposed extension as a partially observable decision problem with external state, internal computational state, and remaining budget. We have implemented and evaluated a RARRL-style PPO reasoning scheduler as a baseline. Across five training seeds, it achieved 70.8 $\pm$ 0.5\% and 85.5 $\pm$ 0.5\% success on two reconstructed abstract tasks, exceeding threshold baselines while using fewer reasoning invocations. An adapted AI2-THOR evaluation with real LLM calls produced 62\% and 57\% success for two frozen schedulers, compared with 57\% for the threshold baseline. A four-cup pilot additionally separates information availability from decision errors. These results motivate further controlled tests of information acquisition and reasoning, rather than establish the full proposed policy.
\end{abstract}

\section{Introduction and Research Objective}
A household robot faces both incomplete observations and ambiguous human goals. A robot instructed to move a particular box may be unable to see its label; changing viewpoint provides evidence that reasoning alone cannot recover. If the labels are visible but the intended box is unclear, asking the human can resolve a different uncertainty. If the evidence is sufficient but not adequately processed, further reasoning may help. Thus, $\mathrm{ACT}_{\rm info}(q)$ can represent either physical inspection or human clarification; the current experiments implement inspection.

ReAct interleaves reasoning and interaction~\cite{yao2023react}, while RARRL selects reasoning configurations under resource costs~\cite{liu2026rarrl}. Two course readings sharpen the remaining question. Hu et al. combine action-induced uncertainty reduction with shielding-aware dual control~\cite{hu2024dual}; KnowNo uses conformal prediction to trigger human assistance when an LLM planner is uncertain~\cite{ren2023knowno}. Their decision rules address interactive information gathering and calibrated help, respectively, without jointly selecting paid LLM reasoning levels and information-acquisition actions. Act to Think asks:
\begin{quote}
\emph{When should an embodied agent first acquire information to improve subsequent reasoning, and how should its representation of uncertainty guide that choice?}
\end{quote}
The proposed Type I contribution is a policy that jointly selects task progress, information acquisition, and internal computation by their downstream value under a shared budget. The scope is high-level decisions with a fixed executor. The completed scheduler and inspection pilot are foundations; they do not yet validate the full learned policy.

This connects to \href{https://alliance-ai.cs.jhu.edu/teaching/HCR/}{course topics}: dynamic programming (Lecture 3), runtime filtering (Lectures 5--6), and POMDP planning with hidden human states (Lecture 14). The budget mask applies the filtering principle to resource feasibility, without providing a physical-safety guarantee. Clarification links the formulation to conformal uncertainty and LLM assistance (Lectures 17 and 19).

\section{Proposed Problem Formulation}
\subsection{Decision events, evidence, and internal state}
Let $n$ index high-level decision events, including both reasoning and external execution. The hidden external state $x_n$ includes physical state and, when modeled, latent human intent. This differs from the robot's computational state below. $H_n$ contains external observations, actions, and execution feedback. Define
\begin{equation}
 b_n(x)=P(x_n=x\mid H_n).
\end{equation}
An internal state $c_n$ stores retained evidence, interpretations, and plans; $B_n$ is the remaining resource budget. The proposed policy chooses
\begin{equation}
\begin{split}
u_n\in\{&\mathrm{ACT}_{\rm task}(q),\ \mathrm{ACT}_{\rm info}(q),\\
        &\mathrm{THINK}(\ell,k)\}.
\end{split}
\label{eq:actions}
\end{equation}
Here $q$ is a symbolic goal, such as picking an object, inspecting its label, or asking which object is intended. The fixed executor maps that goal to an external action, $a_n=E(q,o_n,c_n)$. The policy selects goals and modes; the executor's implementation remains fixed. The reasoning role and computation budget are $\ell$ and $k$.

For an external action, the environment and internal state update as
\begin{equation}
\begin{split}
x_{n+1}&\sim P(\cdot\mid x_n,a_n),\quad
o_{n+1}\sim O(\cdot\mid x_{n+1},a_n),\\
H_{n+1}&=H_n\oplus(a_n,o_{n+1}),\\
c_{n+1}&=U(c_n,a_n,o_{n+1}).
\end{split}
\end{equation}
For THINK in the initial static setting,
\begin{equation}
\begin{aligned}
 c_{n+1}&\sim K_{\ell,k}(\cdot\mid c_n,H_n),\\
 (H_{n+1},x_{n+1})&=(H_n,x_n).
\end{aligned}
\end{equation}
THINK advances the decision-event index but adds no external evidence. A dynamic extension would evolve $x_n$ over its actual duration. Let $C_n^I\geq0$ charge inspection or clarification and $C_n^C\geq0$ charge reasoning, converted into common budget units using fixed rates. Then $d_n=C_n^I+C_n^C$ and $B_{n+1}=B_n-d_n$. Pure task execution has $C_n^I=C_n^C=0$; its time and movement costs enter the task reward and horizon. Thinking has value through bounded computation or a restricted representation; an ideal controller with an exact belief and free computation would not need a charged THINK operation.

\begin{figure}[t]
\centering
\begingroup
\scriptsize
\setlength{\unitlength}{0.95mm}
\begin{picture}(85,61)
\put(0,40){\framebox(32,10){\shortstack{History $H_n\to\widehat b_n$\\State $c_n$; budget $B_n$}}}
\put(48,40){\framebox(32,10){\shortstack{High-level policy\\$\pi_\theta(u_n\mid\widehat b_n,c_n,B_n)$}}}
\put(0,20){\framebox(32,10){\shortstack{Frozen reasoner\\Update $c_n$ from $H_n$}}}
\put(48,20){\framebox(32,10){\shortstack{Fixed executor $E$\\Symbolic goal $q\to a_n$}}}
\put(48,0){\framebox(32,10){Environment}}
\put(32,45){\vector(1,0){16}}
\put(57,40){\line(0,-1){5}}
\put(16,35){\line(1,0){41}}
\put(16,35){\vector(0,-1){5}}
\put(25,36){\makebox(20,3){\footnotesize THINK}}
\put(70,40){\vector(0,-1){10}}
\put(72,33){\makebox(10,3){\footnotesize ACT}}
\put(8,30){\vector(0,1){10}}
\put(64,20){\vector(0,-1){10}}
\put(66,13){\makebox(10,3){$a_n$}}
\put(80,5){\line(1,0){4}}
\put(84,5){\line(0,1){50}}
\put(16,55){\line(1,0){68}}
\put(16,55){\vector(0,-1){5}}
\put(28,56){\makebox(50,3){\footnotesize Observation and execution feedback}}
\end{picture}
\endgroup
\caption{Proposed information and control flow. The current baseline learns reasoning-level selection within a prescribed execution process; selectable information acquisition and the learned belief representation are extensions.}
\label{fig:architecture}
\end{figure}

\subsection{Belief estimation and objective}
A candidate Transformer history encoder and distributional head produce $h_n=f_\phi(H_n)$ and $\widehat b_n=q_\phi(x_n\mid H_n)$. In a finite simulator with exact posteriors, they can be trained using
\begin{equation}
 \mathcal L_{\rm belief}=\E_{H_n}\left[D_{\rm KL}(b_n\Vert\widehat b_n)\right].
\end{equation}
Latent-state cross-entropy is an alternative training signal. Hidden states are not supplied at deployment. The architecture alone does not guarantee a sufficient statistic; approximate information-state theory motivates evaluating reward prediction, controlled prediction, and recursive updates~\cite{subramanian2022ais}.

The policy $\pi_\theta(u_n\mid\widehat b_n,c_n,B_n)$ maximizes
\begin{equation}
 \E_{\pi_\theta}\left[\sum_{n=0}^{N}\gamma^n
 \bigl(r_n^{\rm task}-\lambda_I C_n^I-\lambda_C C_n^C\bigr)\right].
\end{equation}
Here $C_n^I,C_n^C$ are the same charges deducted from $B_n$, and $\lambda_I,\lambda_C\geq0$ convert them into reward penalties. The hard mask admits an operation only if its cost upper bound satisfies $\overline d_n(u_n)\leq B_n$; the soft penalties rank affordable choices by cost-adjusted return. Strict feasibility requires that upper bound to hold, an assumption revisited in Sec.~IV-D.

\section{Implemented Baseline and Training Framework}
\subsection{Task model and baselines}
The implemented baseline addresses a narrower decision than Eq.~\eqref{eq:actions}. The current execution step is prescribed, and the scheduler chooses $k\in\{0,1,2\}$. Level 0 executes directly; levels 1 and 2 represent planning and planning with verification. A nonzero-level step includes paid reasoning followed by execution. The input concatenates task features, remaining budget, recent level--outcome pairs, and pooled attempts and failures; it is not the proposed belief representation.

The reconstructed task-decomposition environment progresses through navigation, inspection, pickup, carrying, and placement; structured navigation contains a longer navigation segment. Some failures terminate the episode, while others permit retry. Reasoning changes failure probability through
\begin{equation}
 p_{\rm fail}(s,k)=p_{\rm base}(s)\bigl(1-\eta_k u(s)\bigr).
\label{eq:failure}
\end{equation}
Base failure rates are 0.06/0.36 for easy/hard navigation regions and 0.12, 0.15, and 0.10 for inspection, pickup, and placement. The uncertainty signal is $u(s)=0.85u_{\rm local}(s)+0.15f_{\rm recent}$, where $u_{\rm local}$ is configured by stage, candidate ambiguity, and local graph structure, and $f_{\rm recent}$ is the recent failure fraction. It is a simulator-defined feature, not a learned posterior. Effectiveness values are $\eta_{0,1,2}=(0,0.5,0.9)$, and reasoning costs are 0, 5--10, and 11--20 units, respectively, against an initial budget of 50--100 units.

The Threshold baseline requests level 2 when $u(s)\ge\tau$ and otherwise level 0, falling back to an affordable level when necessary. Its threshold is selected from $\{0.25,0.40,0.55,0.70,0.85\}$ using validation episode return. Fixed-interval baselines are similarly tuned on validation data. There are no LLM calls during abstract training: it tests allocation under the specified cost--effectiveness model.

\subsection{PPO actor--critic learning}
Only the scheduler's actor and critic are trained with PPO~\cite{schulman2017ppo}. They are separate three-hidden-layer ReLU MLPs; the principal results use width 64. Eight environments each supply 64 transitions per batch. GAE estimates advantages, PPO clipping updates the actor, and squared return-prediction error trains the critic. The configuration uses four update epochs, learning rate $3\times10^{-4}$, discount $0.99$, GAE parameter $0.95$, and clipping parameter $0.2$.

The abstract reward is
\begin{equation}
 r_t=\mathbf{1}_{\rm completion}-0.5\mathbf{1}_{\rm failure}
 -\frac{1+C_t^{\rm reasoning}}{400}.
\end{equation}
Here $t$ indexes the baseline's combined reasoning/execution steps, rather than the separate events of Sec.~II. Episodes last at most 50 such steps. A mask excludes levels whose maximum abstract cost exceeds the remaining budget. Training uses 1,048,576 environment steps per task and seed on CPU. Neither LLM nor executor parameters are updated. The DP solver is an evaluation reference, not a source of training demonstrations.

\section{Progress and Preliminary Results}
\subsection{Reconstructed abstract tasks}
The implementation includes a simulator, NumPy PPO, budget masks, tuned baselines, checkpoints, and DP evaluation. Each task has 1,000 training, 200 validation, and 200 test scenarios. Each test scenario is replayed 20 times with matched random draws across policies. PPO selects the most probable feasible level at evaluation. Table~\ref{tab:abstract} reports five training seeds, 100--104.

\begin{table}[t]
\centering
\caption{Abstract evaluation. SR = replay success rate (\%); RI = mean reasoning invocations. PPO SR is mean $\pm$ sample standard deviation across five training seeds. All policies obey the budget. The DP row uses the same replay metric, not its analytical expectation.}
\label{tab:abstract}
\small
\begin{tabular}{@{}lrrrr@{}}
\toprule
&\multicolumn{2}{c}{Task decomposition}&\multicolumn{2}{c}{Navigation}\\
Policy&SR&RI&SR&RI\\
\midrule
Never think&43.6&0.00&53.2&0.00\\
Always think&55.5&4.69&71.3&4.69\\
Fixed interval&53.2&2.58&64.2&2.03\\
Threshold&60.9&4.57&78.0&4.38\\
PPO&$70.8\!\pm\!0.5$&3.46&$85.5\!\pm\!0.5$&2.62\\
DP (success)&74.4&4.64&86.7&4.63\\
\bottomrule
\end{tabular}
\end{table}

PPO improves success over Threshold by approximately 9.9 and 7.4 percentage points, respectively, using unrounded means, while reducing reasoning use. The success-maximizing DP oracle has exact expected success rates of 74.50\% and 85.59\% under the reconstructed model. Its finite-replay success rates differ because of execution randomness; replay values are not strict analytical upper bounds. These results support learned allocation under hard budgets but do not independently establish LLM reasoning quality, since Eq.~\eqref{eq:failure} prescribes its abstract benefit.

\subsection{Adapted 3D evaluation and calibration}
The 3D pipeline uses AI2-THOR 5.0.0, ALFRED-derived task definitions, a fixed executor, and real GPT-6 Luna calls. At level 0, the executor chooses object and receptacle types by keyword overlap with the instruction, breaking ties by distance; processing requirements use keyword rules, followed by scripted interaction routines. Level 1 invokes the planner with effort \texttt{none}; it remains a paid LLM call with generated output. Level 2 invokes the planner with effort \texttt{low}, rather than a separate verifier. Candidates come from simulator metadata, so this is high-level orchestration rather than standard end-to-end visual ALFRED evaluation.

Calibration measures each stage's error rates and token costs for the executor and both LLM levels. These become stage-specific base risks and effectiveness parameters in an abstract simulator aligned to the five-stage 3D pipeline. Effectiveness matches measured error ratios at a reference uncertainty and can be negative: processing errors are 16\% for the executor, 27\% for level 1, and 12\% for level 2. One resource unit is approximately 36.1 tokens; calibrated level-1/2 costs span 5--14 and 5--15 units. PPO is trained in that calibrated simulator, then frozen for 3D evaluation; it is not fine-tuned on 3D returns.

The same 100 tasks are evaluated under all seven policies in Table~\ref{tab:3d}, totaling 700 episodes: 40 simple placements and 20 each requiring heating, cooling, or cleaning. The two available evaluated PPO checkpoints, seeds 100 and 101, are both reported; this is a two-checkpoint transfer study, not a five-seed estimate.

\begin{table}[t]
\centering
\caption{All seven 3D policies, 100 tasks each. Calls and tokens are episode means. ``Always'' requests the named level when allowed by the runtime mask; unbudgeted low disables that mask.}
\label{tab:3d}
\small
\begin{tabular}{@{}lrrr@{}}
\toprule
Policy&Success (\%)&Calls&Tokens\\
\midrule
Never think&49&0.00&0\\
Always none&45&5.50&2175\\
Always low&57&5.37&2269\\
Threshold&57&3.78&1345\\
Frozen PPO, seed 100&62&3.21&1325\\
Frozen PPO, seed 101&57&3.62&1347\\
Unbudgeted low&57&5.80&2463\\
\bottomrule
\end{tabular}
\end{table}

Both frozen schedulers use fewer calls than Threshold. Seed 100 improves success over Never think by 13 percentage points, with a paired task-bootstrap 95\% interval of $[4,22]$ points, but its five-point advantage over Threshold is absent in seed 101. PPO uses approximately 41--42\% fewer tokens than budgeted Always low. These are preliminary results: the processing-calibration and 3D evaluation sets share 42 of their 100 task IDs, so they do not constitute a strictly independent held-out test. A disjoint evaluation is required for stronger generalization claims.

\subsection{Controlled information--reasoning pilot}
An additional development pilot uses four cups with hidden star/color labels and a prior conditioned on at least one star: 240 equally likely worlds. Exact evaluation across 100 layouts gives an oracle-guided inspection policy perfect accuracy with 1.73 expected inspections, versus four for full inspection at cost 0.08. This is an ideal-computation reference.

The real-LLM pilot uses 12 paired tasks, six per family, with identical histories, prompts, schemas, and output caps for effort \texttt{none}/\texttt{low}. With all labels observed, both solve 6/6 simple tasks, but only 4/6 and 5/6 relational tasks, respectively. Oracle-guided sensing followed by \texttt{low} solves 12/12 using 23 total inspections; \texttt{none} solves 11/12. These small-sample results distinguish missing evidence from processing errors and motivate testing observation--computation trade-offs. They establish neither a learned observer nor an advantage from withholding information.

\subsection{Scope of the evidence}
The reconstruction changes task dynamics, budget handling, network width, the LLM backend, and runtime reasoning configurations relative to RARRL. Physical execution failures constrain the 3D success ceiling. Actual calls can exceed reserved tokens: PPO budget overruns occur in 1\% and 3\% of episodes for the two seeds, so runtime masking lacks the abstract bounded-cost guarantee. The 3D inspection sequence remains prescribed. The cup pilot uses oracle-guided sensing, not a learned belief estimator or information-acquisition policy; human clarification is not evaluated.

\section{Remaining Work and Evaluation Plan}
The finite hidden-state cup task is the primary final-project environment; renewed AI2-THOR transfer is optional. It isolates information--computation choices from physical-execution failures while retaining the observe--reason--execute interface. Exact beliefs and reference action values support diagnosis; sensor noise and human clarification remain extensions.

Next, expand inspection costs and scale from 4 to 8, 16, and 32 cups at controlled observation fractions. Compare raw history, a structured known/unknown inventory, and exact label-belief assistance under matched \texttt{none}/\texttt{low} settings. Paired posterior regret, success, tokens, and latency will distinguish representation assistance from extra computation. The larger experiment is prepared but has no results yet; exact belief assistance is not a learned encoder.

If these checks reveal a useful trade-off, train a small policy over task goals, inspections, and THINK, comparing it with fixed schedules, Threshold, and the existing ACT/THINK baseline. Start with transparent history features and exact-belief references; add a learned belief encoder only after identifying a representation bottleneck. All deployed baselines receive the same observation tools and resource limits. Success, cost-adjusted return, inspection count, reasoning calls, and token use are primary metrics; belief error and action regret diagnose failures. The final evaluation will separate development from held-out tasks and use multiple training seeds. The minimum deliverable is a reproducible information--computation study; neither a learned belief encoder nor renewed 3D transfer is a prerequisite for reporting a valid negative result.

\section{Conclusion}
Act to Think studies how an embodied agent should coordinate new evidence, internal computation, and task execution. The completed baseline improves success and reduces reasoning use in reconstructed hard-budget tasks. Adapted 3D evaluation shows a seed-dependent advantage over a strong heuristic. The cup pilot exposes decision errors even with full evidence, motivating controlled scaling and independent evaluation before learning the joint policy.

\bibliographystyle{plainnat}
\bibliography{references}
\end{document}

